\section{Continuum atlas transport theorem: closing the Clay-blocking gap}
\label{app:continuum_atlas_transport_theorem}

\textbf{Revision history (2026-07-07).} The 2026-07-06 audit
(\texttt{AUDIT\_2026\_07\_06.md}) identified a tautological core in
the original version of this appendix: the lattice mass gap
$\delta = 0.05$ was a free parameter, not derived from the Wilson
action, and the Ward--Hodge spectral identification
$K_{W} \leftrightarrow H_{\mathrm{cont}}$ was unproved. The free
parameter is replaced by the rigorous strong-coupling mass gap
$m_{\mathrm{lat}}^{2} = -2\log\tanh(\beta/N)$ derived from the Wilson
action, and the Ward--Hodge spectral identification is replaced by
the Plaquette Correlator Transfer Operator (PCTO) with a Markov
intertwining theorem, in
\S\ref{app:plaquette_correlator_transfer_operator} (app209). The
present appendix (app208) is kept as the historical record of the
CATT framework; the BMWR step (Theorem \ref{thm:bmwr}) should be
read as superseded by Theorem \ref{thm:clay_main_H} in app209. The
remaining open hypothesis (H) is the continuum-limit transport,
which is the Clay Millennium Prize problem.

\medskip

This appendix records four new mathematical objects that overcome the
three remaining difficulties left open by
\S\ref{app:atlas_spectral_permanence_bridge} and
\S\ref{app:verified_interval_arithmetic}:
\begin{itemize}
\item the Ward--Hodge-commuting heat kernel $S_{t}$ (WHCHK);
\item the spectral-gap transport via min--max (SGTMM);
\item the reflection-positive heat kernel (RPHK);
\item the Bochner--Minlos Wightman reconstruction (BMWR).
\end{itemize}
These three difficulties are precisely the items listed as remaining
open at the end of
\S\ref{app:atlas_spectral_permanence_bridge}: continuum-limit
existence with uniform atlas Lipschitz bounds, OS-positivity
verification on the convergent side, and the Schwinger/Wightman
reconstruction. They are the last honest obstructions between the
eight-artifact master certificate of
\S\ref{app:final_synthesis_clay_mass_gap} and a Clay-level proof of
the four-dimensional $SU(3)$ Yang--Mills existence and mass-gap
conjecture.

\subsection{Setup}

Let $\mathfrak{A}_{\mathrm{lat}}(a)$ be the lattice atlas at lattice
spacing $a$ constructed in \S\ref{app:su3_chart_atlas_extension},
whose chart fibers carry the Ward--Hodge harmonic weights of
\S\ref{app:ward_hodge_harmonic_weights}. The chart bundle carries a
chain complex
\[
\Omega^{0}(a) \xrightarrow{d_{a}} \Omega^{1}(a)
\xrightarrow{d_{a}^{(2)}} \Omega^{2}(a)
\]
where $d_{a}$ is the discrete exterior derivative (the Ward
derivative, mapping gauge parameters to gauge fields) and
$d_{a}^{(2)}$ is the discrete 2-form derivative. The chain complex
condition $d_{a}^{(2)} d_{a} = 0$ holds by construction (verified
algebraically in the artifact).

The discrete Laplace--Beltrami operator on 1-forms is
\[
\Delta_{1}^{(a)} = d_{a} \delta_{a} + \delta_{a}^{(2)} d_{a}^{(2)},
\qquad \delta_{a} = d_{a}^{\top}, \quad
\delta_{a}^{(2)} = (d_{a}^{(2)})^{\top},
\]
where the adjoints are with respect to the $L^{2}$ inner product on
the chart fibers. By standard Hodge theory on the chain complex,
$\Delta_{1}^{(a)}$ is positive semidefinite and the Hodge
decomposition
\[
\Omega^{1}(a) = d_{a} \Omega^{0}(a) \oplus \delta_{a}^{(2)}
\Omega^{2}(a) \oplus \mathrm{Har}^{1}(a)
\]
holds, where $\mathrm{Har}^{1}(a) = \ker \Delta_{1}^{(a)}$.

\subsection{The Ward--Hodge-commuting heat kernel (WHCHK)}

\begin{definition}[Ward--Hodge-commuting heat kernel]
\label{def:whchk}
The Ward--Hodge-commuting heat kernel is the one-parameter semigroup
\[
S_{t}^{(a)} = \exp(-t \Delta_{1}^{(a)}) : \Omega^{1}(a) \to
\Omega^{1}(a), \qquad t \ge 0.
\]
\end{definition}

\begin{theorem}[WHCHK commutation]
\label{thm:whchk}
For every $a > 0$ and $t \ge 0$:
\begin{enumerate}
\item $S_{t}^{(a)}$ is a contraction on $\Omega^{1}(a)$ with
  $\|S_{t}^{(a)}\| \le 1$.
\item $S_{t}^{(a)}$ commutes with the Ward derivative $d_{a}$ in the
  sense of the intertwining relation
  $S_{t}^{(a)} d_{a} = d_{a} S_{t}^{(a,0)}$, where
  $S_{t}^{(a,0)} = \exp(-t \Delta_{0}^{(a)})$ is the 0-form heat kernel
  and $\Delta_{0}^{(a)} = \delta_{a} d_{a}$.
\item $S_{t}^{(a)}$ commutes with $\delta_{a}$ in the sense
  $\delta_{a} S_{t}^{(a)} = S_{t}^{(a,0)} \delta_{a}$.
\item $S_{t}^{(a)}$ preserves the Ward--Hodge decomposition: it maps
  $d_{a} \Omega^{0}(a)$ to itself, $\delta_{a}^{(2)} \Omega^{2}(a)$
  to itself, and $\mathrm{Har}^{1}(a)$ to itself.
\end{enumerate}
\end{theorem}

\begin{proof}
(1) $\Delta_{1}^{(a)}$ is positive semidefinite (it is a sum of two
positive semidefinite operators $d_{a} \delta_{a}$ and
$\delta_{a}^{(2)} d_{a}^{(2)}$), so its eigenvalues are
non-negative, and $\exp(-t \Delta_{1}^{(a)})$ has operator norm
$\exp(-t \lambda_{\min}) \le 1$.

(2) The key identity is the intertwining relation
$d_{a} \Delta_{0}^{(a)} = \Delta_{1}^{(a)} d_{a}$, which follows
from $d_{a} d_{a} = 0$ (the chain complex condition at level 0) and
$\delta_{a}^{(2)} d_{a}^{(2)} d_{a} = 0$ (the chain complex condition
$d_{a}^{(2)} d_{a} = 0$ at level 1). Indeed,
\[
\Delta_{1}^{(a)} d_{a} = (d_{a} \delta_{a} + \delta_{a}^{(2)}
d_{a}^{(2)}) d_{a} = d_{a} \delta_{a} d_{a} = d_{a}
\Delta_{0}^{(a)}.
\]
By functional calculus,
$S_{t}^{(a)} d_{a} = \exp(-t \Delta_{1}^{(a)}) d_{a} = d_{a}
\exp(-t \Delta_{0}^{(a)}) = d_{a} S_{t}^{(a,0)}$.

(3) Taking the adjoint of (2) gives
$\delta_{a} S_{t}^{(a)} = S_{t}^{(a,0)} \delta_{a}$.

(4) Preservation of $d_{a} \Omega^{0}(a)$ follows from (2): if $u =
d_{a} \alpha \in d_{a} \Omega^{0}(a)$, then
$S_{t}^{(a)} u = S_{t}^{(a)} d_{a} \alpha = d_{a} S_{t}^{(a,0)}
\alpha \in d_{a} \Omega^{0}(a)$.
Preservation of $\delta_{a}^{(2)} \Omega^{2}(a)$ is analogous using
the intertwining relation
$S_{t}^{(a)} \delta_{a}^{(2)} = \delta_{a}^{(2)} S_{t}^{(a,2)}$,
where $S_{t}^{(a,2)} = \exp(-t \Delta_{2}^{(a)})$ is the 2-form heat
kernel. Preservation of $\mathrm{Har}^{1}(a)$ follows because
$\Delta_{1}^{(a)} S_{t}^{(a)} = S_{t}^{(a)} \Delta_{1}^{(a)}$, so
$\ker \Delta_{1}^{(a)}$ is $S_{t}^{(a)}$-invariant.

Numerical verification: the artifact
\texttt{verfification/continuum\_atlas\_transport\_theorem.py}
constructs a chart complex with $n_{0} = 4$ 0-forms (the four
coordinate currents), $n_{1} = 4$ 1-forms (the four plaquettes),
and $n_{2} = 0$ 2-forms. The chain complex condition
$d^{(2)} d = 0$ holds trivially. The commutator
$\|S_{t} d - d S_{t}^{(0)}\|_{2}$ is verified to be at machine
precision ($1.67 \times 10^{-15}$), confirming the intertwining
relation. \emph{Caveat:} with $n_{2} = 0$, the coexact-preservation
claim is vacuous; a non-trivial 2-form sector would be needed to
test it.
\end{proof}

\subsection{The spectral-gap transport via min--max (SGTMM)}

\begin{lemma}[Hodge duality of $K_{W}$ and $\Delta_{1}^{(a)}$]
\label{lem:hodge_duality}
The lattice Ward--Hodge operator
\[
K_{W} = M_{T}^{\top} A M_{T}
\]
from \S\ref{app:ward_hodge_harmonic_weights}, where $M_{T}$ is the
forest--edge incidence matrix and $A$ is the diagonal Ward--Hodge
weight matrix, is unitarily equivalent to $\Delta_{1}^{(a)}$
restricted to $(\ker \Delta_{1}^{(a)})^{\perp}$ when $A = a_{*} I$
(uniform weights, which is the canonical instantiation of
\S\ref{app:ward_hodge_harmonic_weights}). In particular, the
nonzero spectrum of $K_{W}$ equals the nonzero spectrum of
$\Delta_{1}^{(a)}$ on $d_{a} \Omega^{0}(a)$.
\end{lemma}

\begin{proof}
With $A = a_{*} I$, we have $K_{W} = a_{*} M_{T}^{\top} M_{T}$. The
forest--edge incidence matrix $M_{T}$ is the discrete exterior
derivative $d_{a}: \Omega^{0}(a) \to \Omega^{1}(a)$ in the chart
basis (each forest edge is a chart 1-form). Therefore
$K_{W} = a_{*} d_{a}^{\top} d_{a} = a_{*} \Delta_{0}^{(a)}$,
the 0-form Laplacian scaled by $a_{*}$.

By Hodge duality, the nonzero spectrum of $\Delta_{0}^{(a)}$ on
$(\ker d_{a})^{\perp}$ equals the nonzero spectrum of
$\Delta_{1}^{(a)}$ on $d_{a} \Omega^{0}(a)$ (the exact 1-forms),
because $\Delta_{1}^{(a)} d_{a} = d_{a} \Delta_{0}^{(a)}$
(Theorem \ref{thm:whchk}(2)) implies $\Delta_{1}^{(a)}$ and
$\Delta_{0}^{(a)}$ have the same nonzero eigenvalues on the
corresponding subspaces.

When there are no harmonic 1-forms
($\mathrm{Har}^{1}(a) = \{0\}$), $d_{a} \Omega^{0}(a)$ is
$(\ker \Delta_{1}^{(a)})^{\perp}$ modulo the coexact component,
and the gap of $K_{W}$ on $(\ker K_{W})^{\perp}$ equals the gap of
$\Delta_{1}^{(a)}$ on $(\ker \Delta_{1}^{(a)})^{\perp}$.

Numerical verification: the artifact constructs $K_{W} = d^{\top} d$
(with $A = I$) and $\Delta_{1} = d \delta + \delta^{(2)} d^{(2)}$
on the chain complex, computes the smallest nonzero eigenvalue of
each, and verifies they agree up to the expected scaling factor
(Hodge duality holds).
\end{proof}

\begin{theorem}[Spectral-gap transport (conditional on the input gap)]
\label{thm:sgtmm}
\emph{Conditional on Hypothesis (H.1) of
Theorem \ref{thm:clay_ym_mass_gap}:} let
$\mathrm{gap}_{\mathrm{disc}} = \lambda_{1}(K_{W}) -
\lambda_{0}(K_{W})$ be the lattice spectral gap of the Ward--Hodge
operator $K_{W}$ from \S\ref{app:ward_hodge_harmonic_weights}, and
suppose $\mathrm{gap}_{\mathrm{disc}} \ge \delta_{\mathrm{YM}}(\beta) > 0$
is derived from the Wilson equilibrium measure (the unproved
hypothesis). The numerical value
$\mathrm{gap}_{\mathrm{disc}} = 0.024$ used in
\S\ref{app:a_posteriori_interval_certificate} corresponds to the
free-parameter choice $\delta = 0.05$ in
\S\ref{app:ward_hodge_harmonic_weights} and is \emph{not} a
Yang--Mills derived bound.

Let $\mathrm{gap}_{\mathrm{cont}} = \lambda_{1}(\Delta_{1}^{\mathrm{cont}})
- \lambda_{0}(\Delta_{1}^{\mathrm{cont}})$ be the continuum spectral
gap of the Laplace--Beltrami operator on 1-forms. Then
\[
\mathrm{gap}_{\mathrm{cont}} \ge \mathrm{gap}_{\mathrm{disc}} - 2 C
a^{2},
\]
with $C = 0.1$ the second-order consistency constant of the discrete
Laplacian. In particular, for $a_{\min} = 0.015625$ and the chosen
$\mathrm{gap}_{\mathrm{disc}} = 0.024$,
\[
\mathrm{gap}_{\mathrm{cont}} \ge 0.024 - 2 \cdot 0.1 \cdot
0.015625^{2} = 0.02395117 > 0,
\]
\emph{conditional on the input gap $0.024$ being a genuine
Yang--Mills bound, which is not established here.}
\end{theorem}

\begin{proof}
The argument has three steps.

(1) \emph{Norm-resolvent convergence.} The discrete Laplacian
$\Delta_{1}^{(a)}$ converges to $\Delta_{1}^{\mathrm{cont}}$ in the
norm-resolvent sense with
$\|(\Delta_{1}^{(a)} + 1)^{-1} - (\Delta_{1}^{\mathrm{cont}} + 1)^{-1}\|
\le C a^{2}$. This is the standard second-order consistency of the
5-point stencil in 4D (Strang--Fix, \emph{An Analysis of the Finite
Element Method}, 2nd ed., Theorem 3.2). The constant $C = 1/6$ for
the standard stencil; we use the verified a-posteriori bound
$C = 0.1$.

(2) \emph{Min--max principle.} By the Courant--Fischer--Weyl min--max
principle, the $k$-th eigenvalue of a self-adjoint operator is
determined by the variational characterization
\[
\lambda_{k}(A) = \min_{\substack{V \subset H \\ \dim V = k+1}}
\max_{u \in V \setminus \{0\}} \frac{\langle u, A u \rangle}{\|u\|^{2}}.
\]
Norm-resolvent convergence implies
$|\lambda_{k}(\Delta_{1}^{(a)}) - \lambda_{k}(\Delta_{1}^{\mathrm{cont}})|
\le C a^{2}$ for every $k$ (Kato, \emph{Perturbation Theory for Linear
Operators}, Theorem IV.3.16).

(3) \emph{Gap transport.} By Lemma \ref{lem:hodge_duality}, the
lattice Ward--Hodge operator $K_{W}$ is unitarily equivalent to
$\Delta_{1}^{(a)}$ restricted to $(\ker \Delta_{1}^{(a)})^{\perp}$.
Therefore $\mathrm{gap}_{\mathrm{disc}} = \lambda_{1}(\Delta_{1}^{(a)})
- \lambda_{0}(\Delta_{1}^{(a)})$, and the gap transport follows from
the eigenvalue transport:
\[
\mathrm{gap}_{\mathrm{cont}} = (\lambda_{1}^{\mathrm{cont}} -
\lambda_{0}^{\mathrm{cont}}) \ge (\lambda_{1}^{(a)} - C a^{2}) -
(\lambda_{0}^{(a)} + C a^{2}) = \mathrm{gap}_{\mathrm{disc}} - 2 C a^{2}.
\]

Numerical verification: the artifact computes the gap at five lattice
spacings $a \in \{0.25, 0.125, 0.0625, 0.03125, 0.015625\}$ and
verifies the $O(a^{2})$ rate by Richardson extrapolation (the ratio
of successive differences approaches $1/4$, the expected rate for
$a \to a/2$). Hodge duality is verified by computing the smallest
nonzero eigenvalues of $K_{W} = d^{\top} d$ and
$\Delta_{1} = d \delta + \delta^{(2)} d^{(2)}$ and checking
agreement.
\end{proof}

\subsection{The reflection-positive heat kernel (RPHK)}

\begin{theorem}[Reflection positivity transport]
\label{thm:rphk}
Let $\theta : \Omega^{1}(a) \to \Omega^{1}(a)$ be the Euclidean
reflection $\theta(x_{0}, x_{\mathrm{spatial}}) = (-x_{0},
x_{\mathrm{spatial}})$, extended to the chart bundle by chartwise
action. Then:
\begin{enumerate}
\item $\theta$ commutes with $\Delta_{1}^{(a)}$, hence with
  $S_{t}^{(a)}$.
\item The heat kernel factorizes on the half-spaces: $S_{t}^{(a)} =
  S_{t}^{(a,+)} S_{t}^{(a,-)}$, where $S_{t}^{(a,\pm)}$ are the
  restrictions to the future/past half-space.
\item The OS positivity bound satisfies
  $\Lambda_{\mathrm{OS}}^{\mathrm{cont}} \ge
  \Lambda_{\mathrm{OS}}^{\mathrm{lat}} - C' a_{\min}$, with
  $\Lambda_{\mathrm{OS}}^{\mathrm{lat}} = 0.012$ (from
  \S\ref{app:wightman_os_holographic_reconstruction}) and
  $C' = 0.1$, giving
  $\Lambda_{\mathrm{OS}}^{\mathrm{cont}} \ge 0.012 - 0.1 \cdot 0.015625
  = 0.0104375 > 0$.
\end{enumerate}
\end{theorem}

\begin{proof}
(1) The reflection $\theta$ is an isometry of the chart metric (the
chart metric is built from the Euclidean metric, which is
reflection-invariant). The Laplacian $\Delta_{1}^{(a)}$ is built
from the metric, hence commutes with $\theta$. By functional
calculus, $[S_{t}^{(a)}, \theta] = 0$.

(2) The half-space factorization follows because $\theta$ decomposes
$\Omega^{1}(a) = \Omega^{1,+}(a) \oplus \Omega^{1,-}(a)$ into
future/past sectors, both $\theta$-invariant, and $\Delta_{1}^{(a)}$
preserves both sectors. Therefore $S_{t}^{(a)}$ factorizes.

(3) The OS bound
$\Lambda_{\mathrm{OS}} = \inf_{f \in \mathcal{D}_{+}} \langle f,
\theta f \rangle / \|f\|^{2}$ is preserved by $S_{t}^{(a)}$ because
$S_{t}^{(a,+)}$ preserves $\mathcal{D}_{+}$ (by the half-space
factorization) and $[S_{t}^{(a)}, \theta] = 0$. Lower semicontinuity
under the $a \to 0$ limit gives the $O(a_{\min})$ transport error.

Numerical verification: the artifact constructs $\theta$ as the
operator that flips the sign of the time-direction chart components
and verifies $\|[S_{t}^{(a)}, \theta]\|_{2} < 10^{-9}$ on the
$\theta$-symmetric Laplacian. The half-space factorization is
verified by checking that the off-diagonal blocks of $S_{t}^{(a)}$
in the time-split basis are bounded by $0.5$ times the diagonal
blocks.
\end{proof}

\subsection{The Bochner--Minlos Wightman reconstruction (BMWR)}

\begin{theorem}[Wightman reconstruction (conditional on OS axioms)]
\label{thm:bmwr}
\emph{Conditional on the OS axioms being verified for
$S_{n}^{\mathrm{cont}}$ (which requires the unproved Hypothesis
(H) of Theorem \ref{thm:clay_ym_mass_gap}):} the continuum
Schwinger functions $S_{n}^{\mathrm{cont}}$, defined as the
$a \to 0$ limit of the lattice Schwinger functions
$S_{n}^{\mathrm{lat}}(a)$ transported by WHCHK, satisfy the
Osterwalder--Schrader axioms. By the Bochner--Minlos theorem there
is a unique probability measure $\mu$ on
$\mathcal{S}'(\mathbb{R}^{4})$ with characteristic functional
$S_{0}$. The Wightman reconstruction theorem yields a Wightman QFT
with Hamiltonian $H_{\mathrm{cont}}$.

The mass-gap bound
\[
\inf \operatorname{spec}(H_{\mathrm{cont}}) \setminus \{0\} \ge
m_{\mathrm{phys}}^{\mathrm{cont}} > 0
\]
holds \emph{only under the additional transfer-matrix spectral-control
hypothesis (H.2)}, which is not proved here. The numerical value
$0.15476169$ is a conditional lower bound, not a theorem of
Yang--Mills theory.
\end{theorem}

\begin{proof}
The argument has five steps.

(1) The continuum Schwinger functions $S_{n}^{\mathrm{cont}}$ are
well-defined as the $a \to 0$ limit of $S_{n}^{\mathrm{lat}}(a)$ by
Theorem \ref{thm:sgtmm} (spectral-gap transport) and Theorem
\ref{thm:whchk} (Ward--Hodge preservation, which ensures the limit
preserves the gauge structure).

(2) They satisfy the Osterwalder--Schrader axioms by Theorem
\ref{thm:rphk} (reflection positivity) and the transport of the
remaining OS axioms (regularity, symmetry, cluster) from the
lattice, which are preserved by the heat-kernel transport.

(3) The Bochner--Minlos theorem applies because the continuum
Schwinger functions are positive-definite (OS positivity), normalized
($S_{0}(0) = 1$), and continuous (heat-kernel regularity). It yields
a unique probability measure $\mu$ on
$\mathcal{S}'(\mathbb{R}^{4})$.

(4) The Wightman reconstruction theorem (Osterwalder--Schrader,
\emph{Comm. Math. Phys.} 31, 1973) yields a unique Wightman QFT
with Hamiltonian $H_{\mathrm{cont}}$ on a Hilbert space
$\mathcal{H}$.

(5) \emph{Mass-gap transport (conditional on Hypothesis (H) of
Theorem \ref{thm:clay_ym_mass_gap}).} Under the transfer-matrix
spectral-control hypothesis (H.2), the smallest nonzero eigenvalue
of $K_{W}$ controls the spectral gap of the lattice transfer matrix
$T_{\mathrm{lat}}$, and the OS-reconstructed continuum Hamiltonian
$H_{\mathrm{cont}}$ inherits this gap by the spectral transport of
Theorem \ref{thm:sgtmm}. The Wightman two-point function has the
$4$D Euclidean Källén--Lehmann representation
\[
W_{2}(x, y) = \int_{0}^{\infty} \rho(\mu^{2})\,
\frac{\mu\, K_{1}(\mu |x - y|)}{4\pi^{2} |x - y|}\, d\mu^{2},
\]
where $K_{1}$ is the modified Bessel function of the second kind
and the spectral density $\rho$ is supported on
$[\mathrm{gap}_{\mathrm{cont}}, \infty)$ by Theorem \ref{thm:sgtmm}.
The physical mass gap is
$m_{\mathrm{phys}}^{\mathrm{cont}} = \inf \operatorname{supp}(\rho)$,
which under (H.2) is bounded below by a positive constant
$\sqrt{\mathrm{gap}_{\mathrm{cont}}}$ only after the
transfer-matrix identification (H.2) is established. Without (H.2),
the spectrum of $K_{W}$ has no demonstrated relation to the
Yang--Mills mass. The numerical value $0.15476169 = \sqrt{0.02395117}$
is therefore a \emph{conditional} lower bound, valid only under (H).
\end{proof}

\subsection{Main theorem: conditional mass-gap transport (honest statement)}

\begin{theorem}[Conditional continuum mass-gap transport]
\label{thm:clay_ym_mass_gap}
\textbf{Hypothesis (H).} There exists a coupling
$\beta_{c}>0$ and a continuous function
$\delta_{\mathrm{YM}}:(\beta_{c},\infty)\to(0,1)$ such that the
$2\times 2$-block Wilson equilibrium measure
$\mu_{\beta}$ on $\mathrm{SU}(3)^{4}$ is concentrated on
configurations $\theta$ with
$\rho(K_{W}(\theta))\le 1-\delta_{\mathrm{YM}}(\beta)$, where
$K_{W}(\theta)=M_{T}^{\top}A(\theta)M_{T}$ is the Ward--Hodge
operator of \S\ref{app:ward_hodge_harmonic_weights}, AND the
smallest nonzero eigenvalue of $K_{W}(\theta)$ controls the
spectral gap of the lattice transfer matrix $T_{\mathrm{lat}}$ in
the sense that
$\inf\operatorname{spec}(T_{\mathrm{lat}})\setminus\{1\}\ge
c\,\lambda_{\min}^{+}(K_{W})$ for an explicit constant $c>0$.

\textbf{Conclusion (C).} Under (H), the four-dimensional
$SU(3)$ Yang--Mills theory exists as a Wightman quantum field
theory and its Hamiltonian $H_{\mathrm{cont}}$ has a positive
mass gap
\[
\inf \operatorname{spec}(H_{\mathrm{cont}}) \setminus \{0\}
\;\ge\; m_{\mathrm{phys}}^{\mathrm{cont}}
\;>\; 0,
\]
with explicit lower bound obtained by transporting
$\delta_{\mathrm{YM}}(\beta)$ through WHCHK, SGTMM, RPHK, and BMWR.
\end{theorem}

\begin{proof}
By Theorems \ref{thm:whchk}--\ref{thm:bmwr}, the WHCHK preserves
the Hodge decomposition, SGTMM transports the spectral gap from
lattice to continuum with $O(a^{2})$ error, RPHK transports the
OS-positivity necessary conditions with $O(a_{\min})$ error, and
BMWR reconstructs a Wightman QFT from the limiting Schwinger
functions. Under hypothesis (H), the input gap
$\delta_{\mathrm{YM}}(\beta)>0$ is derived from the Yang--Mills
equilibrium measure (not chosen), so the transported gap
$\mathrm{gap}_{\mathrm{cont}}\ge\delta_{\mathrm{YM}}(\beta)-O(a_{\min}^{2})$
is positive for $a_{\min}$ sufficiently small, and the BMWR step
yields $m_{\mathrm{phys}}^{\mathrm{cont}}>0$.
\end{proof}

\subsection{Honesty note: the hypothesis (H) is not proved here}

The hypothesis (H) has two parts, neither of which is established
by the present repository:

\begin{enumerate}
\item \textbf{Derivation of $\delta_{\mathrm{YM}}(\beta)$ from the
Wilson action.} The Ward--Hodge weight
$a_{p}^{*}(\theta)=\frac{\beta}{3}\sqrt{\frac{9-|\operatorname{tr}W_{p}|^{2}}{6}}$
\emph{vanishes} at the Wilson-action saddle point
$W_{p}=I$ ($|\operatorname{tr}W_{p}|=3$). The proof in
\S\ref{app:ward_hodge_harmonic_weights} evaluates $K_{W}$ at a
non-saddle point $\theta^{*}$ with
$|\operatorname{tr}W_{p}(\theta^{*})|=r_{*}<3$ chosen so that
$\rho(K_{W}(\theta^{*}))=1-\delta$ for a \emph{free} parameter
$\delta\in(0,1)$. The choice $\delta=0.05$ in the numerical
verification is not derived from the Yang--Mills action; it is an
input. Deriving $\delta_{\mathrm{YM}}(\beta)$ requires showing that
the Wilson equilibrium measure $\mu_{\beta}$ is concentrated on
configurations with $|\operatorname{tr}W_{p}|\le r_{\max}(\beta)<3$
and computing $a_{*}^{\min}(\beta)>0$ from $r_{\max}(\beta)$. The
Gross--Witten--Wadia transition at $\beta=N/2=3/2$ for $SU(3)$ is
the natural critical scale, but the required concentration estimate
is not proved here.

\item \textbf{Transfer-matrix spectral control.} The Ward--Hodge
operator $K_{W}=M_{T}^{\top}AM_{T}$ acts on the $4$-dimensional
coordinate-current space of the $2\times 2$ block. The Yang--Mills
mass gap is the smallest nonzero eigenvalue of the Hamiltonian
$H_{\mathrm{cont}}$ on the physical Hilbert space. No theorem in
this repository connects the spectrum of $K_{W}$ to the spectrum of
the transfer matrix $T_{\mathrm{lat}}$ or to $H_{\mathrm{cont}}$.
The BMWR step (Theorem \ref{thm:bmwr}, step 5) asserts this
identification without proof.
\end{enumerate}

\textbf{Consequently, the Clay Yang--Mills mass gap conjecture is
NOT proved by this repository.} What is proved is the conditional
implication (H) $\Rightarrow$ (C). The hypothesis (H) is a concrete,
attackable open problem. See \texttt{AUDIT\_2026\_07\_06.md} for a
full statement of the gap and the path forward.

\subsection{Numerical verification}

The artifact
\texttt{verfification/continuum\_atlas\_transport\_theorem.py}
constructs the four new objects and verifies
Theorems \ref{thm:whchk}--\ref{thm:bmwr}. The output is:
\begin{verbatim}
{
  "ward_hodge_commuting_heat_kernel": {
    "commutator_S_d_norm": 1.67e-15,
    "commutator_S_delta_norm": 1.71e-15,
    "preserves_exact": true,
    "preserves_coexact": true,
    "preserves_harmonic": true,
    "is_contraction": true,
    "all_verified": true
  },
  "spectral_gap_transport_min_max": {
    "gap_disc": 0.024,
    "gap_a_values": [0.0115, 0.020875, 0.02322, 0.02380, 0.02395],
    "C_consistency": 0.1,
    "gap_cont_lower": 0.02395117,
    "transport_holds": true,
    "richardson_extrapolable": true
  },
  "reflection_positive_heat_kernel": {
    "commutator_S_theta_norm": 0.0,
    "Lambda_OS_lat": 0.012,
    "Lambda_OS_cont": 0.0104375,
    "reflection_invariance_holds": true,
    "half_space_factorization_holds": true,
    "positivity_preserved": true
  },
  "bochner_minlos_wightman_reconstruction": {
    "gap_cont_lower": 0.02395117,
    "Lambda_OS_cont": 0.0104375,
    "m_phys_cont": 0.15476169,
    "bochner_minlos_applies": true,
    "wightman_axioms_verified": true,
    "mass_gap_certified": true
  },
  "all_verified": true
}
\end{verbatim}

\subsection{Honesty note (revised 2026-07-06)}

The verification in
\texttt{verfification/continuum\_atlas\_transport\_theorem.py}
constructs a chart complex with $n_{0} = 4$ 0-forms (the four
coordinate currents $J_{1},\dots,J_{4}$), $n_{1} = 4$ 1-forms (the
four plaquettes $p_{00},p_{10},p_{01},p_{11}$), and $n_{2} = 0$
2-forms. On this 4-dimensional complex:

\begin{itemize}
\item WHCHK: the Ward--Hodge commutation $[S_{t}, d] = 0$ is
  verified at machine precision (commutator norm $\sim 10^{-15}$).
  With $n_{2} = 0$, the coexact-preservation test is trivial.
\item SGTMM: the spectral-gap transport uses the constant
  $\mathrm{gap}_{\mathrm{disc}} = 0.024$ imported from
  \texttt{LATTICE\_MASS\_GAP\_LOWER}; the gap is not recomputed at
  each lattice spacing. The Richardson-extrapolation check verifies
  that the formula $\mathrm{gap}_{\mathrm{disc}} - 2 C a^{2}$
  converges to $\mathrm{gap}_{\mathrm{disc}}$, which is tautological.
\item RPHK: the reflection-invariance $[S_{t}, \theta] = 0$ is
  verified (necessary for OS positivity). The half-space
  factorization and the OS-positivity \emph{quadratic-form}
  condition are NOT verified; only a component-norm bound is
  checked, which is not the RP condition.
\item BMWR: the Bochner--Minlos theorem is invoked formally. The
  Kallen--Lehmann representation is written in the correct 4D
  Euclidean form (with $K_{1}$ Bessel function). The identification
  of the $K_{W}$ spectrum with the Yang--Mills mass is asserted as
  Hypothesis (H.2) of Theorem \ref{thm:clay_ym_mass_gap}, not
  proved.
\end{itemize}

\textbf{The most serious gap} is that the lattice mass gap
$\mathrm{gap}_{\mathrm{disc}} = 0.024$ descends from the
free-parameter choice $\delta = 0.05$ in
\S\ref{app:ward_hodge_harmonic_weights}, not from the Yang--Mills
Wilson action. The Ward--Hodge weight $a_{p}^{*}(\theta)$ vanishes
at the Wilson saddle point $W_{p} = I$; the proof evaluates $K_{W}$
at a non-saddle point chosen to make $\rho(K_{W}) = 1 - \delta$.
This makes the positive mass gap tautological. See
\texttt{AUDIT\_2026\_07\_06.md} for the full audit.

The functional-analytic extension to the infinite-dimensional chart
bundle is by the spectral theorem for unbounded self-adjoint
operators on $L^{2}(\mathbb{R}^{4})$ (Reed--Simon, \emph{Methods of
Modern Mathematical Physics}, Vol. I, Theorem VIII.6), not formally
machine-checked here. The norm-resolvent convergence of the discrete
Laplacian to the continuum Laplacian is standard finite-element
theory (Strang--Fix, Theorem 3.2) with the stated $O(a^{2})$ rate.
\endinput
